\section{Prime and divisor queries: ready-to-use code}
\textbf{Select by bounds.} For many prime queries with maximum endpoint
at most $10^7$, build \texttt{PrimeTable} once. It provides primality,
the ordered prime list, and inclusive range counts. Its arrays cost about
$5(V+1)$ bytes plus the prime list, roughly 53 MB at $V=10^7$ with
4-byte \texttt{int}. Check the contest memory limit.
For a narrow interval near a large number, use the segmented routine below.
It still needs primes through $\sqrt R$; a short interval alone does not
make arbitrary 64-bit endpoints cheap.
\sourceimport{finals-number}{Prelude.h}

\subsection{Is x prime? Count or list primes in [L,R]}
0 and 1 are not prime. Cross off multiples beginning at $p^2$: smaller
multiples already have a smaller prime factor. Build
$P[x]=\#\{p\le x:p\text{ prime}\}$. Then the answer on inclusive
$[L,R]$ is $P[R]-P[L-1]$, with the $L=0$ boundary handled separately.
\sourceimport{finals-number}{PrimeQueries.h}
\begin{lstlisting}
using namespace finals_number;
PrimeTable t(1000000); // build ONCE, outside query loop
bool yes = t.is_prime(97); // true
int count = t.count(10, 30); // 6
vector<int> p = t.list(10, 30);
// p = {11, 13, 17, 19, 23, 29}
\end{lstlisting}
Build $O(V\log\log V)$; primality and range count $O(1)$; listing takes
$O(\log\pi(V)+k)$ for $k$ returned primes. Every queried endpoint must be
within the built table. Rebuilding for each query wastes the preprocessing.

\textbf{Complete query program.} Copy Prelude and PrimeQueries above, then
add this main. Input: maximum endpoint $V$, query count $Q$, then $Q$ pairs
$L,R$. Output: one count per line. All queries obey $0\le L\le R\le V$.
\begin{lstlisting}
#include <iostream>
int main() {
    std::ios::sync_with_stdio(false);
    std::cin.tie(nullptr);
    int V, Q;
    std::cin >> V >> Q;
    finals_number::PrimeTable table(V);
    while (Q--) {
        int L, R;
        std::cin >> L >> R;
        std::cout << table.count(L, R) << '\n';
    }
}
\end{lstlisting}
For input \texttt{30 3} followed by \texttt{0 1}, \texttt{2 2},
\texttt{10 30}, the output lines are 0, 1, 6. A compilable version is
\texttt{handbooks/examples/prime\_range\_queries.cpp} in the repository.

\subsection{Count prime elements in an array interval}
This is a different question: count occurrences in array positions, not
prime integers in a numeric interval. Duplicates must count repeatedly.
After building \texttt{t} through the largest nonnegative array value:
\begin{lstlisting}
std::vector<int> pref(a.size() + 1, 0);
for (int i = 0; i < (int)a.size(); ++i)
    pref[i + 1] = pref[i] + (a[i] >= 2 &&
                            t.is_prime(a[i]));
// Array positions [l,r), not numeric values [l,r].
int answer = pref[r] - pref[l];
\end{lstlisting}
For $a=[2,4,2,7]$, the whole array contains three prime occurrences,
although it contains only two distinct prime values. Positions $[1,4)$
contain two. With point updates, store the prime indicator in a Fenwick
tree and update by new-indicator minus old-indicator. The prime table must
cover every possible updated value too.

\subsection{High endpoints: segmented prime enumeration}
\texttt{segmented\_primes(L,R)} marks just the inclusive interval and
returns its primes in sorted order. For each base prime $p\le\sqrt R$,
the first marked multiple is
$\max(p^2,\lceil L/p\rceil p)$; using only the rounded multiple would
incorrectly mark $p$ itself when it lies in the interval.
\sourceimport{finals-number}{SegmentedPrimes.h}
\begin{lstlisting}
auto p = finals_number::segmented_primes(10, 30);
auto count = p.size(); // 6
// Large endpoint, small interval also allowed:
auto high = finals_number::segmented_primes(
    1000000000000LL - 1000, 1000000000000LL);
\end{lstlisting}
This implementation enforces $0\le L\le R\le10^{12}$ and width
$W=R-L+1\le10^7$. It needs $O(\sqrt R+W)$ space, including output,
and scans the base-prime list per call. Many high-interval queries may need
shared base preprocessing. To count only and reduce output memory, count
surviving flags instead of storing them; the supplied API returns the list.
A huge width is not supported just because $R$ is within the endpoint bound.
Perfect-square endpoints, 0, 1, and single-number intervals are tested.

\subsection{How many divisors? What is their sum?}
If many positive queries are bounded by $V$, add each possible divisor $d$
to all multiples of $d$. This builds the divisor count and divisor sum for
every value in $O(V\log V)$ time and $O(V)$ memory. Query in $O(1)$.
\sourceimport{finals-number}{DivisorTable.h}
\begin{lstlisting}
finals_number::DivisorTable d(1000000);
int count = d.count[12]; // 6
long long sum = d.sum[12]; // 28
long long proper_sum = d.sum[12] - 12; // 16
bool perfect = (d.sum[6] - 6 == 6); // true
// d.count[1] == 1, d.sum[1] == 1
\end{lstlisting}
Only positive query values are valid; zero has no finite divisor count.
At $V=2\cdot10^6$, the two arrays need about 24 MB. If only counts are
needed, omit the sums array and its update. For a few large integers, use
factorization or divisor enumeration instead of allocating up to their value.
The repository also includes runnable \texttt{divisor\_queries.cpp} and
\texttt{segmented\_primes.cpp} examples beside the prime query program.
